\chapter{Parámetros S, redes y estabilidad}\label{cap:sparams}

\section{Por qué parámetros S}

Los parámetros Z, Y, H o ABCD se miden con circuitos abiertos o cortocircuitos en los puertos. En RF eso es imposible: un abierto
radia y tiene capacidad parásita, un corto tiene inductancia, y un transistor puede oscilar con esas terminaciones. Los parámetros S
(\emph{scattering}, de dispersión) se miden con todos los puertos terminados en $Z_0$ y relacionan \textbf{ondas}, que es justo lo que mide un VNA.

\section{Ondas de potencia}
\begin{definicion}{ondas $a$ y $b$}
En cada puerto $i$ con tensión $V_i$ y corriente $I_i$ (entrante), y una referencia real $Z_0$:
\begin{equation}
  a_i = \frac{V_i + Z_0 I_i}{2\sqrt{Z_0}}, \qquad b_i = \frac{V_i - Z_0 I_i}{2\sqrt{Z_0}} .
\end{equation}
\begin{donde}
\sym{a_i}{onda incidente en el puerto $i$ (entra en la red)}{$\sqrt{\mathrm{W}}$}
\sym{b_i}{onda reflejada o emergente del puerto $i$ (sale de la red)}{$\sqrt{\mathrm{W}}$}
\sym{\abs{a_i}^2}{potencia incidente (con valores eficaces)}{\si{\watt}}
\end{donde}
\end{definicion}
Comparando con el capítulo~\ref{cap:lineas}: $a_i = V_i^+/\sqrt{Z_0}$ y $b_i = V_i^-/\sqrt{Z_0}$. Dividir por $\sqrt{Z_0}$ hace que
$\abs{a}^2$ sea directamente una potencia.

\section{La matriz S}
\begin{definicion}{matriz de dispersión}
\begin{equation}
  \mathbf{b} = \mathbf{S}\,\mathbf{a}, \qquad
  \begin{pmatrix} b_1 \\ b_2 \end{pmatrix} = \begin{pmatrix} S_{11} & S_{12} \\ S_{21} & S_{22} \end{pmatrix}\begin{pmatrix} a_1 \\ a_2 \end{pmatrix},
  \qquad S_{ij} = \left.\frac{b_i}{a_j}\right|_{a_k = 0,\ k\ne j}.
\end{equation}
\begin{donde}
\sym{S_{ij}}{respuesta que sale por el puerto $i$ cuando se excita sólo el puerto $j$ y los demás están terminados en $Z_0$}{---}
\sym{S_{11}, S_{22}}{coeficientes de reflexión de entrada y de salida (con el otro puerto adaptado)}{---}
\sym{S_{21}}{transmisión directa: ganancia o pérdida de inserción}{---}
\sym{S_{12}}{transmisión inversa: aislamiento}{---}
\end{donde}
\end{definicion}
Notación del primer subíndice: \emph{dónde se mide}; del segundo: \emph{dónde se excita}. $S_{21}$ se lee «del 1 al 2».

\subsection{Magnitudes habituales}
\begin{align}
  IL &= -20\log_{10}\abs{S_{21}} \quad\text{(pérdida de inserción, positiva)}\\
  G &= 20\log_{10}\abs{S_{21}} \quad\text{(ganancia, para amplificadores)}\\
  \tau_g &= -\frac{\dd\phi_{21}}{\dd\omega} = -\frac{1}{2\pi}\frac{\dd\phi_{21}}{\dd f}
\end{align}
\begin{donde}
\sym{\phi_{21}}{fase de $S_{21}$, desenrollada (sin saltos de $360^\circ$)}{rad}
\sym{\tau_g}{retardo de grupo: tiempo que tarda la envolvente de una señal en atravesar la red}{\si{\second}}
\end{donde}
\begin{intuicion}[: retardo de grupo]
Si la fase cae linealmente con la frecuencia, todas las componentes de una señal llegan con el mismo retardo y la forma se conserva.
Un retardo de grupo variable en la banda distorsiona señales de banda ancha (es la razón por la que se especifica la «planitud» del
retardo de grupo en filtros).
\end{intuicion}

\subsection{Propiedades de redes físicas}
\begin{itemize}
\item \textbf{Recíproca} (sin ferritas ni elementos activos): $\mathbf S = \mathbf S^T$, es decir, $S_{12} = S_{21}$.
\item \textbf{Sin pérdidas}: $\mathbf S^H\mathbf S = \mathbf I$ (matriz unitaria). Para dos puertos: $\abs{S_{11}}^2 + \abs{S_{21}}^2 = 1$.
\item \textbf{Pasiva}: los autovalores de $\mathbf S^H\mathbf S$ son $\le 1$; ninguna combinación de entradas devuelve más potencia de la que recibe.
\end{itemize}
La cantidad $1 - \abs{S_{11}}^2 - \abs{S_{21}}^2$ es la fracción de potencia disipada (o radiada) en la red.

\section{Conversión entre representaciones}
Con $Z_0$ real igual en todos los puertos y $\mathbf I$ la matriz identidad:
\begin{equation}
  \mathbf S = (\mathbf Z - Z_0\mathbf I)(\mathbf Z + Z_0\mathbf I)^{-1}, \qquad \mathbf Z = Z_0(\mathbf I + \mathbf S)(\mathbf I - \mathbf S)^{-1}.
\end{equation}
Es la generalización matricial de $\Gamma = (Z - Z_0)/(Z + Z_0)$. Para admitancias, con $Y_0 = 1/Z_0$:
$\mathbf S = (Y_0\mathbf I - \mathbf Y)(Y_0\mathbf I + \mathbf Y)^{-1}$. El motor convierte H y G a S con las fórmulas cerradas de Frickey
(válidas para dos puertos) y comprueba todas las conversiones contra scikit-rf en los tests.

\subsection{Matriz ABCD y cascadas}
\begin{definicion}{parámetros de transmisión}
\begin{equation}
  \begin{pmatrix} V_1 \\ I_1 \end{pmatrix} = \begin{pmatrix} A & B \\ C & D \end{pmatrix}\begin{pmatrix} V_2 \\ -I_2 \end{pmatrix}.
\end{equation}
\begin{donde}
\sym{A}{relación de tensiones con salida en abierto, $V_1/V_2$}{---}
\sym{B}{impedancia de transferencia con salida en corto}{\si{\ohm}}
\sym{C}{admitancia de transferencia con salida en abierto}{\si{\siemens}}
\sym{D}{relación de corrientes con salida en corto}{---}
\end{donde}
\end{definicion}
Su gran ventaja: \textbf{la red total de una cascada es el producto de las matrices}, en orden desde la entrada:
\begin{equation}
  \mathbf T_{\text{total}} = \mathbf T_1\,\mathbf T_2\cdots\mathbf T_n .
\end{equation}
Los bloques básicos son:
\begin{equation}
  \underbrace{\begin{pmatrix} 1 & Z \\ 0 & 1\end{pmatrix}}_{\text{serie } Z}, \qquad
  \underbrace{\begin{pmatrix} 1 & 0 \\ Y & 1\end{pmatrix}}_{\text{paralelo } Y}, \qquad
  \underbrace{\begin{pmatrix} \cosh\gamma\ell & Z_0\sinh\gamma\ell \\ \sinh(\gamma\ell)/Z_0 & \cosh\gamma\ell\end{pmatrix}}_{\text{línea}} .
\end{equation}
Y la impedancia de entrada con una carga $Z_L$ al final es
\begin{equation}
  Z_{in} = \frac{A Z_L + B}{C Z_L + D}.
\end{equation}
\begin{enlaapp}
Así funciona por dentro el \app{Circuit Builder} y la verificación del \app{Matching Designer}: cada elemento (serie, paralelo, línea, stub)
es una matriz ABCD por frecuencia, se multiplican todas y se aplica la fórmula de $Z_{in}$. Los puntos intermedios de la cascada dan los
arcos dibujados en la carta de Smith.
\end{enlaapp}

\section{Reflexión de entrada con carga arbitraria}
Si el puerto 2 no está terminado en $Z_0$ sino en una carga con coeficiente $\Gamma_L$, entonces $a_2 = \Gamma_L b_2$ y despejando:
\begin{equation}
  \Gamma_{in} = S_{11} + \frac{S_{12}S_{21}\Gamma_L}{1 - S_{22}\Gamma_L}, \qquad
  \Gamma_{out} = S_{22} + \frac{S_{12}S_{21}\Gamma_S}{1 - S_{11}\Gamma_S}.
\end{equation}
\begin{donde}
\sym{\Gamma_L, \Gamma_S}{coeficientes de reflexión de la carga y de la fuente, referidos a $Z_0$}{---}
\sym{\Gamma_{in}, \Gamma_{out}}{reflexión vista en la entrada y en la salida del dispositivo}{---}
\end{donde}
El segundo término es la «realimentación» a través de $S_{12}$: sólo si $S_{12} = 0$ (red unilateral) la entrada no depende de la carga.

\section{Estabilidad de amplificadores}

Un amplificador es \textbf{inestable} si para alguna carga o fuente pasiva ($\abs{\Gamma_L}, \abs{\Gamma_S} < 1$) resulta
$\abs{\Gamma_{in}} > 1$ o $\abs{\Gamma_{out}} > 1$: presenta resistencia negativa y puede oscilar.

\begin{definicion}{criterios de estabilidad incondicional}
\begin{align}
  \Delta &= S_{11}S_{22} - S_{12}S_{21}\\
  K &= \frac{1 - \abs{S_{11}}^2 - \abs{S_{22}}^2 + \abs{\Delta}^2}{2\,\abs{S_{12}S_{21}}} > 1 \quad\text{y}\quad \abs\Delta < 1\\
  \mu &= \frac{1 - \abs{S_{11}}^2}{\abs{S_{22} - \Delta\conj{S_{11}}} + \abs{S_{12}S_{21}}} > 1
\end{align}
\begin{donde}
\sym{\Delta}{determinante de la matriz S}{---}
\sym{K}{factor de Rollett; necesita además la condición $\abs\Delta<1$}{---}
\sym{\mu}{factor de Edwards--Sinsky: condición necesaria y suficiente por sí sola; además, cuanto mayor, más margen}{---}
\end{donde}
\end{definicion}
\begin{intuicion}[: qué mide $\mu$]
$\mu$ es la distancia desde el centro de la carta de $\Gamma_L$ hasta la región inestable más cercana. $\mu > 1$ significa que toda la carta
(todas las cargas pasivas) queda en zona estable. Por eso la aplicación lo usa como indicador principal.
\end{intuicion}

\subsection{Círculos de estabilidad}
El lugar de las cargas $\Gamma_L$ que dan exactamente $\abs{\Gamma_{in}} = 1$ es una circunferencia en el plano $\Gamma_L$:
\begin{equation}
  C_L = \frac{\conj{\left(S_{22} - \Delta\conj{S_{11}}\right)}}{\abs{S_{22}}^2 - \abs\Delta^2}, \qquad
  r_L = \abs{\frac{S_{12}S_{21}}{\abs{S_{22}}^2 - \abs\Delta^2}} .
\end{equation}
\begin{donde}
\sym{C_L, r_L}{centro (complejo) y radio del círculo de estabilidad de salida}{---}
\end{donde}
Para saber qué lado es estable basta un punto de prueba: el centro de la carta ($\Gamma_L = 0$) da $\abs{\Gamma_{in}} = \abs{S_{11}}$; si es
menor que 1, el lado del círculo que contiene el origen es el estable. Análogamente para la fuente ($C_S$, $r_S$, intercambiando 1 y 2).

\begin{ejemplo}{transistor de microondas (datos de ejemplo)}
$S_{11} = 0.894\angle{-60.6^\circ}$, $S_{12} = 0.020\angle 62.4^\circ$, $S_{21} = 3.122\angle 123.6^\circ$, $S_{22} = 0.781\angle{-27.6^\circ}$:
\[
  K = \StbK, \qquad \abs\Delta = \StbD, \qquad \mu = \StbMu .
\]
Como $\mu < 1$ (y $K<1$), el transistor es \textbf{potencialmente inestable}: hay cargas pasivas que lo hacen oscilar, y la red de
adaptación debe evitar las regiones marcadas por los círculos de estabilidad.
\end{ejemplo}

\subsection{Ganancias y círculos de ganancia}
\begin{itemize}
\item \textbf{Ganancia máxima estable} (si $K<1$): $MSG = \abs{S_{21}}/\abs{S_{12}}$.
\item \textbf{Ganancia máxima disponible} (si $K>1$): $MAG = \frac{\abs{S_{21}}}{\abs{S_{12}}}\left(K - \sqrt{K^2 - 1}\right)$, que se alcanza con
adaptación conjugada simultánea de entrada y salida.
\item \textbf{Círculos de ganancia}: lugar de $\Gamma_L$ con ganancia de potencia $G_p$ constante (o de $\Gamma_S$ con ganancia disponible $G_a$
constante). Permiten cambiar ganancia por estabilidad, ruido o ancho de banda.
\end{itemize}

\section{Parámetros en modo mixto}

En líneas diferenciales (USB, Ethernet, LVDS, antenas balanceadas) interesa la señal \emph{diferencia} entre dos conductores y la
\emph{común} a ambos. Para un par de puertos $(p, n)$:
\begin{equation}
  a_d = \frac{a_p - a_n}{\sqrt2}, \qquad a_c = \frac{a_p + a_n}{\sqrt2}.
\end{equation}
Agrupando en una matriz ortogonal $\mathbf M$ (cumple $\mathbf M^{-1} = \mathbf M^T$), los parámetros en modo mixto son
\begin{equation}
  \mathbf S_{mm} = \mathbf M\,\mathbf S\,\mathbf M^{T} = \begin{pmatrix}\mathbf S_{dd} & \mathbf S_{dc}\\ \mathbf S_{cd} & \mathbf S_{cc}\end{pmatrix}.
\end{equation}
\begin{donde}
\sym{\mathbf S_{dd}}{respuesta diferencial a excitación diferencial (lo que interesa)}{---}
\sym{\mathbf S_{cc}}{respuesta en modo común a excitación en modo común}{---}
\sym{\mathbf S_{cd}}{conversión de diferencial a común: genera radiación (EMI)}{---}
\sym{\mathbf S_{dc}}{conversión de común a diferencial: susceptibilidad al ruido}{---}
\end{donde}
El factor $1/\sqrt2$ conserva la potencia: la impedancia diferencial de referencia es $2Z_0$ y la común $Z_0/2$.

\begin{enlaapp}
El módulo \app{S-Parameters} muestra módulo, fase, retardo de grupo, VSWR, conversiones a Z, Y, ABCD, H y G, estabilidad ($K$, $\Delta$,
$\mu$, MSG/MAG) en función de la frecuencia, círculos de estabilidad y de ganancia en la carta, y modo mixto con emparejamiento de puertos consecutivos
(1--2, 3--4\ldots), verificado contra la función equivalente de scikit-rf.
\end{enlaapp}
